\section{An Empirical Programme}\label{sec:empirical}

The constructions above would be of limited interest if they could not be tested, and the central proposal, that a book can teach its own script, is exactly the kind of statement that is easy to find plausible and hard to confirm. This section specifies what the tests would be. It reports no results and is not a pilot. Its purpose is to state the hypotheses, conditions and measures in enough detail that a result, including a negative one, could be interpreted, and that the same design could be run by others with their own material.

\subsection{Hypotheses}

The first group of hypotheses concerns whether acquisition happens at all. \textbf{H1} says that readers of a drifting edition recognise the final-stage forms in isolation more accurately than readers who encounter the same forms for the first time, so that the book has taught something; this is the minimal claim, and it distinguishes acquisition from contextual guessing. \textbf{H2} says that this recognition transfers, in that readers of the drifting edition read a fresh passage, not part of the book, set in the final-stage script with higher accuracy and speed than control readers. The second group concerns the comparison of ways of presenting the transition. \textbf{H3} says that a drifting presentation yields better transfer than presenting the final script from the outset, with or without an alphabet table, at equal reading time. \textbf{H4} says that a drift whose schedule is driven by exposure, in the sense of Section~\ref{sec:drift}, yields better transfer for rare letters than a drift whose schedule is uniform in position. \textbf{H5} says that routes matter: for the same endpoints and the same schedule, routes with designed intermediate keyframes preserve comprehension better during the transition than routes that interpolate outlines linearly. The third group concerns durability and the limits of the claim. \textbf{H6} says that recognition survives a delay of days or weeks, with the possibility that the pattern of retention depends on how exposures were spaced \citep{cepeda2006,bjork1994}. \textbf{H7} says that the receptive gains of H1 and H2 are not accompanied by comparable gains in production, which would confirm that the book teaches the first level of knowing a sign and not the highest. None of these is assumed to be true, and the design is intended to produce interpretable results if all are false.

\subsection{Conditions}

The conditions to be compared fall naturally into a small number of arms, each defined by a manifest in the sense of Section~\ref{sec:source}. A \emph{plain} arm presents the text in ordinary letters and serves as a baseline for comprehension and reading speed. A \emph{static} arm presents the final script from the first page without assistance, and a \emph{static with table} arm presents it with an alphabet table. An \emph{abrupt} arm switches from ordinary to final forms at a single point. The drifting arms differ along the dimensions identified earlier: uniform against exposure-driven schedules, linear against keyframe routes, reciprocal pairs with and without detour marks, new-shape endpoints against swapped-shape endpoints, and monotone drifts against cycles. A complete factorial design is far too large, and the programme proceeds in stages, with the earliest experiments comparing a few arms on a single alphabet and later ones varying one dimension at a time from the best-performing arm.

\subsection{Measures}

The primary measures are chosen to correspond to the levels of knowing a sign distinguished in Section~\ref{sec:arb}. Comprehension is measured by questions about the content of the passage, scored independently of the script, and reading speed by words per unit time, which together show whether a presentation preserves the ability to follow an argument, the criterion by which any transition must be judged. Recognition in isolation is measured by presenting single letter forms, drawn from the final alphabet, and recording accuracy and response time, by letter, so that frequency effects can be examined. Transfer is measured by a fresh passage in the final script. Retention is measured by repeating these tests after a delay. Production is measured separately, by asking participants to write a given word in the final script. Subjective ratings of difficulty are collected but are secondary, since readers are poor judges of what they have learned. For the transitions themselves, performance can also be recorded as a function of position in the book, which would show where in a drift comprehension dips, and this connects the empirical record to the margin and novelty-load quantities of Section~\ref{sec:drift}.

\subsection{Design considerations}

Several features of the setting complicate inference. Participants differ in prior exposure: a reader who knows Standard Galactic from games, or who reads Arabic, begins from a different place than one who does not, and these must be measured and balanced. The specific alphabet confounds the effect of presentation, so the substitution should be counterbalanced across participants, with different pairings used for different readers, or else the study estimates the difficulty of one alphabet and not of a method. The same text should not be seen by the same participant in more than one arm, so that memory for content is not confused with learning of script. Participants should be assigned to arms at random, the exposures each letter receives should be recorded so that the dependence on exposure can be analysed directly, and analysis should use models that treat participants and letters as sampled units. The hypotheses, the exclusion rules and the analysis plan should be recorded before data are collected. The effect sizes are not known in advance, so the size of the first studies would be set from pilots, and the design is meant to report intervals and not merely tests of significance.

There is also precedent for the kind of effect sought. Readers asked to read transformed typography show substantial practice effects, and some of the benefit appears to persist long after the practice \citep{kolers1976}. The work on skilled reading offers a framework for what a measure of reading speed and comprehension can show about the processes underlying it \citep{rayner2012}. A careful reading of the history of the Initial Teaching Alphabet, in which a transitional script was designed to give way to the conventional one, would supply cautions about transfer from a special script back to the usual one \citep{downing1967}.

\subsection{The question to be answered}

The programme can be summarised in one question: which routes preserve comprehension while producing recognition that survives outside the particular book that taught it? The two clauses can come apart, and the distinction between them is the reason for distinguishing contextual assistance from acquisition. A book may be perfectly readable throughout its drift because the reader is guessing from context, and may leave nothing behind. A book may teach the forms but at a cost to comprehension so great that no one finishes it. The target is a route on which both conditions hold, and the empirical programme exists to find out whether such a route exists, for which alphabets, and for which readers. If it does not, that is a result, and the essay's claim that the arbitrariness of the sign is separable from the constraints on its acquisition is unaffected by it, since that claim is precisely about how much of the difficulty lies in the assignment and how much in the structure of the shapes and the reader.
